\section{Exponent ledger and passage beyond the Poisson face}
\label{sec:exponent-passage}

Write
\begin{equation}
 Q=X^{\beta+o(1)},\qquad J=X^{1-\beta+o(1)},
 \qquad S=X^{s+o(1)},\qquad T=X^{t+o(1)},
 \label{eq:asym-exponent-main}
\end{equation}
with \(0<s\le t\).  On a denominator and conductor cell put
\begin{equation}
 Y=X^{y+o(1)},\qquad F_i=X^{f_i+o(1)},
 \qquad k_i=\max\{f_i,\beta-f_i\}.
 \label{eq:asym-cell-exponents}
\end{equation}
The rectangular support gives
\begin{equation}
 0\le y\le s+t,\qquad 0\le f_1\le s,\qquad 0\le f_2\le t,
 \label{eq:asym-cell-ranges}
\end{equation}
and \(\cK_Q(F_i)=X^{k_i+o(1)}\).

The natural scale \(XQ=Q^2J\) has exponent \(1+\beta\).  The norm and
direct-Gauss routes in \eqref{eq:asymmetric-cell-bound} have exponents
\[
 \beta+\frac{1-\beta+y+k_1+k_2}{2},
 \qquad
 \beta+f_1+f_2+\frac{k_1+k_2}{2},
\]
respectively.  Their saving exponents are therefore
\begin{equation}
 \eta_A=\frac{\beta+1-y-k_1-k_2}{2},
 \qquad
 \eta_B=1-f_1-f_2-\frac{k_1+k_2}{2}.
 \label{eq:asym-etaAB}
\end{equation}

\begin{proposition}[Asymmetric three-cell conductor ledger]
\label{prop:asymmetric-ledger}
Every cell in \eqref{eq:asym-cell-ranges} satisfies
\begin{equation}
 \max\{\eta_A,\eta_B\}\ge
 \eta_{\rm asym}(\beta;s,t),
 \label{eq:asym-ledger-bound}
\end{equation}
where
\begin{equation}
 \boxed{
 \eta_{\rm asym}(\beta;s,t)
 =\min\left\{
 1-\frac{3\beta}{2},
 \frac{\beta+1-2s-2t}{2},
 \frac{3-\beta-s-5t}{4}
 \right\}.}
 \label{eq:eta-asymmetric}
\end{equation}
\end{proposition}

\begin{proof}
Decreasing \(y\) improves \(\eta_A\), so use \(y=s+t\).

If \(f_1,f_2\le\beta/2\), then
\(k_1+k_2=2\beta-f_1-f_2\), and
\[
 \eta_B=1-\beta-\frac{f_1+f_2}{2}
 \ge1-\frac{3\beta}{2}.
\]
If \(f_1,f_2\ge\beta/2\), then
\(k_1+k_2=f_1+f_2\le s+t\), whence
\[
 \eta_A\ge\frac{\beta+1-2s-2t}{2}.
\]

In a mixed cell write
\(x\le\beta/2\le z\) and \(d=z-x\ge0\).  Then
\(k_1+k_2=\beta+d\), so
\begin{equation}
 \eta_A\ge\frac{1-s-t-d}{2}.
 \label{eq:asym-mixed-A}
\end{equation}
Since \(z\le t\),
\begin{equation}
 \eta_B=1-\frac\beta2-2z+\frac d2
 \ge1-\frac\beta2-2t+\frac d2.
 \label{eq:asym-mixed-B}
\end{equation}
The maximum is at least the average of the last two lower bounds,
namely \((3-\beta-s-5t)/4\).  Combining the three cases proves the
claim.
\end{proof}

\subsection{The mass-sensitive mixed-increment theorem}

\begin{theorem}[One-sided calibrated-increment estimate]
\label{thm:calibrated-increment-estimate}
Assume the row and target-primitivity interfaces of
\cref{sec:interface-prefixes}, together with the divisor-independent
off-diagonal mask, \eqref{eq:no-wrap}, and
\eqref{eq:Fourier-tail-interface}.  The short-side condition
\eqref{eq:short-side} is not required for this calibrated formulation.
Suppose
\(\eta_{\rm asym}(\beta;s,t)>0\) with a fixed margin and
\(0<\kappa<s\).  Then for every fixed \(A>0\) and every sufficiently
small fixed \(\epsilon>0\), there is a fixed \(\eta>0\) such that
\begin{align}
 |\cI_{S,T}|
 \ll{}&XQX^{-\eta}
 \notag\\
 &+J\cZ_1\left\{
 (\log X)^{-A}+X^{-s+\kappa+o(1)}+L^{-1}X^\epsilon
 \right\}.
 \label{eq:mass-sensitive-main-theorem}
\end{align}
Here \(\eta\) may be any sufficiently small fixed number below the
conductor margin and the principal-mean margin \(1-\beta\), after the
auxiliary \(\epsilon\)'s are chosen.  The transpose packet obeys the
same estimate.
\end{theorem}

\begin{proof}
Use \cref{prop:joint-minus-drift}.  Reassemble the joint \(r\ne0\)
discrepancy through \eqref{eq:asymmetric-cell-reassembly}.  Apply
\cref{prop:asymmetric-ledger} and
\eqref{eq:asymmetric-cell-bound}; the fixed margin absorbs all dyadic
and divisor multiplicities.  The global principal cell is bounded by
\eqref{eq:asymmetric-principal-mean}; under
\eqref{eq:row-pair-mass} it is
\(Q^2X^\epsilon=XQX^{-(1-\beta)+\epsilon}\).  The physical zero mode
is \eqref{eq:mass-sensitive-zero}, and the drift is
\eqref{eq:drift-packet-bound}.  The inherited Fourier tail is made
smaller than the displayed fixed-power term by choosing the smooth
truncation order after its polynomial mass.  This proves
\eqref{eq:mass-sensitive-main-theorem}.
\end{proof}

\begin{corollary}[Two rigorous closure regimes]
\label{cor:two-closure-regimes}
Retain the hypotheses of
\cref{thm:calibrated-increment-estimate}, and suppose
\(L\ge X^{\lambda_0}\) for a fixed \(\lambda_0>0\).
\begin{enumerate}[label=\textup{(\roman*)}]
 \item If the total projective mass satisfies
 \eqref{eq:polylog-projective-mass}, then for every fixed \(B>0\),
 \begin{equation}
  \boxed{\cI_{S,T}\ll_{B,K,F,H}XQ(\log X)^{-B}.}
  \label{eq:polylog-full-closure}
 \end{equation}
 \item Under only the soft mass hypothesis
 \eqref{eq:row-pair-mass}, if \(\widehat F(0)=0\), then some fixed
 \(\eta'>0\) satisfies
 \begin{equation}
  \boxed{\cI_{S,T}\ll XQX^{-\eta'}.}
  \label{eq:mean-zero-full-closure}
 \end{equation}
 after the conventional auxiliary \(X^\epsilon\) losses are chosen
 smaller than the fixed margins.
\end{enumerate}
\end{corollary}

\begin{proof}
In (i), choose the calibration exponent \(A\) after \(B,K\).  The
large-\(g\), drift, conductor, and principal-mean terms have fixed
power savings.  In (ii), the physical zero mode is identically zero by
\cref{cor:mass-sensitive-zero}; every remaining term has a fixed-power
margin.
\end{proof}

\begin{remark}[What the soft theorem does not say]
\label{rem:soft-not-log}
With general \(\widehat F(0)\ne0\) and only
\(\cZ_1\ll Q^2X^\epsilon\), the first term inside braces in
\eqref{eq:mass-sensitive-main-theorem} remains
\(XQX^\epsilon(\log X)^{-A}\).  It is neither asserted to be
\(o(XQ)\) nor relabelled as logarithmic saving.
\end{remark}

\subsection{The general one-sided passage region}

For a genuine crossing of the old Poisson face, require
\begin{equation}
 s<1-\beta<t.
 \label{eq:general-poisson-crossing}
\end{equation}
Positivity of the three entries in \eqref{eq:eta-asymmetric} gives the
explicit window
\begin{equation}
 \boxed{
 1-\beta<t<
 \min\left\{
 \frac{\beta+1}{2}-s,
 \frac{3-\beta-s}{5}
 \right\},}
 \label{eq:general-passage-window}
\end{equation}
together with \(\beta<2/3\) and \(s<1-\beta\), all with fixed
margins.  The interval in \eqref{eq:general-passage-window} is
nonempty precisely when
\begin{equation}
 \max\left\{\frac{1+2s}{3},\frac{2+s}{4}\right\}
 <\beta<\min\left\{\frac23,1-s\right\}.
 \label{eq:general-beta-window}
\end{equation}
There exist \(\beta,t\) satisfying these inequalities if and only if
\begin{equation}
 \boxed{s<\frac25.}
 \label{eq:general-s-frontier}
\end{equation}
These are sufficient faces for the certified three-cell ledger, not a
claim that the two underlying cell bounds are globally optimal.

\subsection{The baseline cutoff
\texorpdfstring{\(S=R\)}{S=R}}

Take
\begin{equation}
 s=\frac12-\delta,\qquad
 t=\frac12-\delta+\xi.
 \label{eq:Rxi-substitution}
\end{equation}
Then \eqref{eq:eta-asymmetric} becomes
\begin{equation}
 \boxed{
 \eta_{\rm asym}(\beta,\delta,\xi)
 =\min\left\{
 1-\frac{3\beta}{2},
 \frac{\beta+4\delta-1}{2}-\xi,
 \frac{6\delta-\beta-5\xi}{4}
 \right\}.}
 \label{eq:eta-asym-special}
\end{equation}
At \(\xi=0\), this agrees with the TPC-22 square ledger; when
\(s=t\), the general formula agrees with the TPC-23 dynamic-square
ledger.

The condition \(T>J\) is
\begin{equation}
 \xi>\frac12+\delta-\beta.
 \label{eq:T-beyond-J}
\end{equation}
Together with positivity of \eqref{eq:eta-asym-special}, this proves
the passage window
\begin{equation}
 \boxed{
 \frac12+\delta-\beta<\xi<
 \min\left\{
 \frac{\beta+4\delta-1}{2},
 \frac{6\delta-\beta}{5}
 \right\},}
 \label{eq:special-passage-window}
\end{equation}
under
\begin{equation}
 \beta<\frac23,\qquad \beta<\frac12+\delta.
 \label{eq:special-beta-margins}
\end{equation}
The latter inequality keeps the baseline leg \(S=R\) strictly below
\(J\).  A simple nonempty subregion is
\begin{equation}
 \boxed{
 \delta>\frac1{10},\qquad
 \frac58-\frac\delta4<\beta<
 \min\left\{\frac23,\frac12+\delta\right\}.}
 \label{eq:special-nonempty-region}
\end{equation}

For the exact rational sample
\begin{equation}
 \delta=\frac3{20},\qquad
 \beta=\frac{31}{50},\qquad
 \xi=\frac1{25},
 \label{eq:rational-passage-sample}
\end{equation}
we have
\[
 R=X^{7/20},\qquad J=X^{19/50},\qquad
 T=X^{39/100}>J,
\]
and the three entries of \eqref{eq:eta-asym-special} are
\[
 \frac7{100},\qquad\frac7{100},\qquad\frac1{50}.
\]
Thus the certified joint nonzero saving is \(1/50\).
